\scanpage{104}
% Archive foliation: 105.
\uline{Dém.} On a (i)$\Longrightarrow$(ii), car (i) implique que $\operatorname{Ker}f_s^j=\operatorname{Im}f_s^{n-j}$, et on conclut par le lemme préliminaire.

\deleted{\unsure{(i)$\Longrightarrow$(ii)}} (ii)$\Longrightarrow$(i), on a $H=\operatorname{Ker}f^j$, considérons $G\xrightarrow{g}H$ induit par $f^{n-j}$ (existe, car $f^n=0$ i.e. $f^j(f^{n-j})=0$). Par hyp. (ii), il est \inserted{fid} plat par fibre. Comme $G$ est plat de prés. finie sur $S$ et $H$ de prés. finie sur $S$, on en conclut que $g$ est \inserted{fid} plat \uline{et} $H$ est plat.

\uline{Cor.} Ceci ne dépend pas du choix de $j$, et implique que \deleted{\unsure{les $\operatorname{Ker}f^j/\operatorname{Ker}f^{j'}$ existent donc}} \inserted{les $\operatorname{Ker}f^j$ sont plats}, \deleted{\unsure{et que}} pour $j\geq j'$ $\operatorname{Ker}f^j/\operatorname{Ker}f^{j'}$ est représentable comme schém. en groupes de type fini, et $f$ dans $\mathcal C$ faisceaux \deleted{\unsure{plats}} fppf \deleted{\unsure{sur}} \inserted{\unsure{plats}} \unsure{vérifie les} conditions de Tate de hauteur $n$.

Cas $G$ fini sur $S$ : rang de Tate en un pt.
\margin{(loc. cte)}
Il y a un seul sous-objet de \unsure{T.A} de hauteur donnée $i\leq n$, de $=$ rang $\rho$.

Kif-kif pour quotients.
